\documentclass[10pt,a4paper]{amsart}
\usepackage{amsmath,amssymb,amsthm,graphicx}
% For commutative diagrams we use Xy-pic
\usepackage[all]{xy}

% We use 2cell for 2-commutative diagrams.
\xyoption{2cell}
% Theorem environments.
%
\theoremstyle{plain}
\newtheorem{theorem}[subsection]{Theorem}
\newtheorem{proposition}[subsection]{Proposition}
\newtheorem{lemma}[subsection]{Lemma}

\theoremstyle{definition}
\newtheorem{definition}[subsection]{Definition}
\newtheorem{example}[subsection]{Example}
\newtheorem{exercise}[subsection]{Exercise}
\newtheorem{situation}[subsection]{Situation}

\theoremstyle{remark}
\newtheorem{remark}[subsection]{Remark}
\newtheorem{remarks}[subsection]{Remarks}

\numberwithin{equation}{subsection}

% Macros
%
\def\lim{\mathop{\mathrm{lim}}\nolimits}
\def\colim{\mathop{\mathrm{colim}}\nolimits}
\def\Spec{\mathop{\mathrm{Spec}}}
\def\Hom{\mathop{\mathrm{Hom}}\nolimits}
\def\Ext{\mathop{\mathrm{Ext}}\nolimits}
\def\SheafHom{\mathop{\mathcal{H}\!\mathit{om}}\nolimits}
\def\SheafExt{\mathop{\mathcal{E}\!\mathit{xt}}\nolimits}
\def\Sch{\mathit{Sch}}
\def\Mor{\mathop{\mathrm{Mor}}\nolimits}
\def\Ob{\mathop{\mathrm{Ob}}\nolimits}
\def\Sh{\mathop{\mathit{Sh}}\nolimits}
\def\NL{\mathop{N\!L}\nolimits}
\def\CH{\mathop{\mathrm{CH}}\nolimits}
\def\proetale{{pro\text{-}\acute{e}tale}}
\def\etale{{\acute{e}tale}}
\def\QCoh{\mathit{QCoh}}
\def\Ker{\mathop{\mathrm{Ker}}}
\def\Im{\mathop{\mathrm{Im}}}
\def\Coker{\mathop{\mathrm{Coker}}}
\def\Coim{\mathop{\mathrm{Coim}}}

% Boxtimes
%
\DeclareMathSymbol{\boxtimes}{\mathbin}{AMSa}{"02}

%
% Macros for moduli stacks/spaces
%
\def\QCohstack{\mathcal{QC}\!\mathit{oh}}
\def\Cohstack{\mathcal{C}\!\mathit{oh}}
\def\Spacesstack{\mathcal{S}\!\mathit{paces}}
\def\Quotfunctor{\mathrm{Quot}}
\def\Hilbfunctor{\mathrm{Hilb}}
\def\Curvesstack{\mathcal{C}\!\mathit{urves}}
\def\Polarizedstack{\mathcal{P}\!\mathit{olarized}}
\def\Complexesstack{\mathcal{C}\!\mathit{omplexes}}
% \Pic is the operator that assigns to X its picard group, usage \Pic(X)
% \Picardstack_{X/B} denotes the Picard stack of X over B
% \Picardfunctor_{X/B} denotes the Picard functor of X over B
\def\Pic{\mathop{\mathrm{Pic}}\nolimits}
\def\Picardstack{\mathcal{P}\!\mathit{ic}}
\def\Picardfunctor{\mathrm{Pic}}
\def\Deformationcategory{\mathcal{D}\!\mathit{ef}}
\newenvironment{reference}{}{}
\newenvironment{slogan}{}{}
\newenvironment{history}{}{}
\usepackage{cleveref}
\crefname{originallemma}{original lemma}{original lemmas}
\Crefname{originallemma}{Original lemma}{Original lemmas}
\makeatletter
\newlabel{curves-lemma-etale-local-structure-nodal-family}{{0CBY}{source}{}{sourcepointer-1}{}}
\newlabel{curves-lemma-etale-local-structure-nodal-family@cref}{{[originallemma][0][]0CBY}{[1][0][]source}{}{}{}}
\makeatother
\crefname{originaldefinition}{original definition}{original definitions}
\Crefname{originaldefinition}{Original definition}{Original definitions}
\makeatletter
\newlabel{more-algebra-definition-extension-discrete-valuation-rings}{{09E4}{source}{}{sourcepointer-2}{}}
\newlabel{more-algebra-definition-extension-discrete-valuation-rings@cref}{{[originaldefinition][0][]09E4}{[1][0][]source}{}{}{}}
\makeatother
\crefname{originallemma}{original lemma}{original lemmas}
\Crefname{originallemma}{Original lemma}{Original lemmas}
\makeatletter
\newlabel{more-algebra-lemma-Dedekind-etale-extension}{{0AP2}{source}{}{sourcepointer-3}{}}
\newlabel{more-algebra-lemma-Dedekind-etale-extension@cref}{{[originallemma][0][]0AP2}{[1][0][]source}{}{}{}}
\makeatother
\crefname{originallemma}{original lemma}{original lemmas}
\Crefname{originallemma}{Original lemma}{Original lemmas}
\makeatletter
\newlabel{morphisms-lemma-smooth-etale-over-affine-space}{{054L}{source}{}{sourcepointer-4}{}}
\newlabel{morphisms-lemma-smooth-etale-over-affine-space@cref}{{[originallemma][0][]054L}{[1][0][]source}{}{}{}}
\makeatother
\title{The étale local structure of an at-worst-nodal scheme over a discrete valuation ring}
\author{The Stacks Project Authors}
\date{stacks-project:models.tex:local-snapshot-20261012; GNU FDL 1.2 or later}
\begin{document}
\maketitle
\noindent\textit{Editorial note.} The counted unit is the Stacks Project Lemma lemma-etale-local-at-worst-nodal at source lines 6106-6130 together with its proof at source lines 6132-6188, a straight adjacent pairing in the chapter file. The chapter is the upstream file models.tex; the pool directory carries the longer name stacks-models and is not used as the chapter name anywhere in this card.\par\medskip
\noindent\textit{Editorial note.} The retained statement contains the author's own commutative diagram drawn with Xy-pic, so the retained preamble keeps the chapter preamble's Xy-pic block at lines 14-18 verbatim. The diagram is part of the statement and is retained without alteration.\par\medskip
\noindent\textit{Editorial note.} This card has no context span: the statement and proof are self-contained apart from the four pointers, and the notions they use are recorded in the standard prerequisites. In particular the notion of being at-worst-nodal is used only in the form recorded there.\par\medskip
\noindent\textit{Editorial note.} The retained proof quotes four results of other chapters, in Curves, More on Algebra and Morphisms; all four are declared as non-load-bearing navigation pointers with their stable Stacks tags looked up in the shared tag index, and none of their statements is reproduced inside the span.\par\medskip
\noindent\textit{Editorial note.} The proof span starts at line 6132, the line of the opening marker itself, which was read from the source rather than inferred from the statement's end.\par\medskip
\noindent\textit{Editorial note.} No citation occurs in the retained spans, so no bibliography entry is supplied.\par\medskip
\noindent\textit{Allowed standard prerequisites (editorial).}\begin{itemize}
\item At-worst-nodal schemes of relative dimension one, étale extensions of discrete valuation rings and the node equation. A scheme of relative dimension one over a discrete valuation ring is at-worst-nodal when its singularities are ordinary double points whose branches are smooth; in a neighbourhood of such a point the local equation can be put in the form of a product of two coordinate functions, possibly with a power of the uniformizer subtracted, after an étale extension of the ring and after an étale modification of the scheme. An étale extension of discrete valuation rings is an extension that is flat and unramified, and Dedekind domains admit étale extensions realising prescribed residue extensions only under a separability condition; a morphism to an affine space which is smooth at a point is étale after shrinking. The author's own lemmas of the chapter and the quoted results of the More on Algebra, Morphisms and Curves chapters are the allowed inputs.. The counted lemma is the local model on which the chapter's resolution of at-worst-nodal schemes by blowups is carried out, and the following lemmas of the section compute with this model.
\end{itemize}\medskip
\ifdefined\hypertarget\hypertarget{sourcepointer-1}{}\fi
\noindent\textit{Source reference.} Original lemma 0CBY: focused-seeds/AGP/sources/stacks-snapshot-20261011/tags.txt, source lines 15488--15488. Allowed input of the Curves chapter quoted by the retained proof for the étale local structure of a nodal family; the pointer is registered so the reference resolves and the cited statement is not reproduced inside the span. The number is the stable Stacks tag of the target, looked up in the shared tag index.\par\medskip
\ifdefined\hypertarget\hypertarget{sourcepointer-2}{}\fi
\noindent\textit{Source reference.} Original definition 09E4: focused-seeds/AGP/sources/stacks-snapshot-20261011/tags.txt, source lines 11887--11887. Allowed input of the More on Algebra chapter quoted by the retained proof for the notion of an extension of discrete valuation rings; the pointer is registered so the reference resolves and the cited definition is not reproduced inside the span. The number is the stable Stacks tag of the target, looked up in the shared tag index.\par\medskip
\ifdefined\hypertarget\hypertarget{sourcepointer-3}{}\fi
\noindent\textit{Source reference.} Original lemma 0AP2: focused-seeds/AGP/sources/stacks-snapshot-20261011/tags.txt, source lines 13460--13460. Allowed input of the More on Algebra chapter quoted by the retained proof to obtain the étale extension of the discrete valuation ring with the required residue field; the pointer is registered so the reference resolves and the cited statement is not reproduced inside the span. The number is the stable Stacks tag of the target, looked up in the shared tag index.\par\medskip
\ifdefined\hypertarget\hypertarget{sourcepointer-4}{}\fi
\noindent\textit{Source reference.} Original lemma 054L: focused-seeds/AGP/sources/stacks-snapshot-20261011/tags.txt, source lines 6644--6644. Allowed input of the Morphisms chapter quoted by the retained proof to recognise the local morphism to the affine model as étale once it is smooth of the right relative dimension; the pointer is registered so the reference resolves and the cited statement is not reproduced inside the span. The number is the stable Stacks tag of the target, looked up in the shared tag index.\par\medskip
\begin{lemma}
\label{lemma-etale-local-at-worst-nodal}
Let $R$ be a discrete valuation ring. Let $X$ be a scheme which is
at-worst-nodal of relative dimension $1$ over $R$.
Let $x \in X$ be a point of the special fibre
of $X$ over $R$. Then there exists a commutative diagram
$$
\xymatrix{
X \ar[d] &
U \ar[r] \ar[d] \ar[l] &
\Spec(A) \ar[dl] \\
\Spec(R) &
\Spec(R') \ar[l]
}
$$
where $R \subset R'$ is an \'etale extension of discrete valuation rings,
the morphism $U \to X$ is \'etale, the morphism $U \to \Spec(A)$ is \'etale,
there is a point $x' \in U$ mapping to $x$, and
$$
A = R'[u, v]/(uv)
\quad\text{or}\quad
A = R'[u, v]/(uv - \pi^n)
$$
where $n \geq 0$ and $\pi \in R'$ is a uniformizer.
\end{lemma}

\begin{proof}
We have already proved this lemma in much greater generality, see
Algebraic Curves, Lemma \ref{curves-lemma-etale-local-structure-nodal-family}.
All we have to do here is to translate the statement
given there into the statement given above.

\medskip\noindent
First, if the morphism $X \to \Spec(R)$ is smooth at $x$,
then we can find an \'etale morphism $U \to \mathbf{A}^1_R = \Spec(R[u])$ for
some affine open neighbourhood $U \subset X$ of $x$. This is
Morphisms, Lemma \ref{morphisms-lemma-smooth-etale-over-affine-space}.
After replacing the coordinate $u$ by $u + 1$ if necessary, we may
assume that $x$ maps to a point in the standard open
$D(u) \subset \mathbf{A}^1_R$. Then $D(u) = \Spec(A)$ with
$A = R[u, v]/(uv - 1)$ and we see that the result is true in this case.

\medskip\noindent
Next, assume that $x$ is a singular point of the fibre. Then we may apply
Algebraic Curves, Lemma \ref{curves-lemma-etale-local-structure-nodal-family}
to get a diagram
$$
\xymatrix{
X \ar[d] &
U \ar[rr] \ar[l] \ar[rd] & &
W \ar[r] \ar[ld] &
\Spec(\mathbf{Z}[u, v, a]/(uv - a)) \ar[d] \\
\Spec(R) & &
V \ar[ll] \ar[rr] & & \Spec(\mathbf{Z}[a])
}
$$
with all the properties mentioned in the statement of the cited lemma.
Let $x' \in U$ be the point mapping to $x$ promised by the lemma.
First we shrink $V$ to an affine neighbourhood of the image of $x'$.
Say $V = \Spec(R')$. Then $R \to R'$ is \'etale. Since $R$ is a
discrete valuation ring, we see that $R'$ is a finite
product of quasi-local Dedekind domains (use
More on Algebra, Lemma \ref{more-algebra-lemma-Dedekind-etale-extension}).
Hence (for example using prime avoidance) we find a standard
open $D(f) \subset V = \Spec(R')$ containing the image of $x'$
such that $R'_f$ is a discrete valuation ring.
Replacing $R'$ by $R'_f$ we reach the situation where
$V = \Spec(R')$ with $R \subset R'$ an \'etale extension of
discrete valuation rings (extensions of discrete valuation rings
are defined in More on Algebra, Definition
\ref{more-algebra-definition-extension-discrete-valuation-rings}).

\medskip\noindent
The morphism $V \to \Spec(\mathbf{Z}[a])$ is determined by
the image $h$ of $a$ in $R'$. Then $W = \Spec(R'[u, v]/(uv - h))$.
Thus the lemma holds with
$A = R'[u, v]/(uv - h)$. If $h = 0$ then we clearly
obtain the first case mentioned in the lemma. If $h \not = 0$
then we may write $h = \epsilon \pi^n$ for some $n \geq 0$
where $\epsilon$ is a unit of $R'$. Changing coordinates
$u_{new} = \epsilon u$ and $v_{new} = v$ we obtain the second
isomorphism type of $A$ listed in the lemma.
\end{proof}

\end{document}
